\documentclass[11pt]{article}
\usepackage[margin=1in]{geometry}
\usepackage{booktabs}
\usepackage{longtable}
\usepackage{array}
\usepackage{xcolor}
\usepackage{fancyhdr}
\usepackage{hyperref}
\usepackage{enumitem}
\usepackage{tcolorbox}
\usepackage{parskip}
\usepackage{titlesec}
\usepackage{gensymb}

\definecolor{accent}{RGB}{44,62,80}
\definecolor{lightgray}{RGB}{245,246,248}
\definecolor{midgray}{RGB}{120,120,120}

\titleformat{\section}{\Large\bfseries\color{accent}}{\thesection}{0.5em}{}
\titleformat{\subsection}{\large\bfseries\color{accent}}{\thesubsection}{0.5em}{}

\pagestyle{fancy}
\fancyhf{}
\renewcommand{\headrulewidth}{0pt}
\renewcommand{\footrulewidth}{0.4pt}
\fancyfoot[C]{\small\color{midgray}Page \thepage\ of \pageref{LastPage}}
\fancyhead[L]{\small\color{midgray}BRISC 2025 -- Edge Deployment Report}
\usepackage{lastpage}

\hypersetup{colorlinks=true, linkcolor=accent, urlcolor=accent, citecolor=accent}

\newcolumntype{L}[1]{>{\raggedright\arraybackslash}p{#1}}

\title{\textbf{\Large Edge Deployment Report}\\[4pt]
\large BRISC 2025 Brain Tumor MRI -- Inference and Edge Deployment}
\author{}
\date{August 27, 2026}

\begin{document}
\maketitle
\thispagestyle{fancy}

\noindent How the BRISC brain tumor classification and segmentation models were taken from
trained checkpoints to working edge deployments on a Raspberry Pi~4 and an NVIDIA
Jetson Nano, including the measured latency, memory, power, and thermal numbers, the
wireless scan-to-report pipeline, and what still needs a human in the loop.

\begin{tcolorbox}[colback=lightgray, colframe=accent, boxrule=0.4pt, arc=1pt]
\small Raspberry Pi 4 (4GB RAM, CPU only) \;\textbullet\; NVIDIA Jetson Nano (4GB RAM,
128-core Maxwell CUDA GPU) \;\textbullet\; Colab T4 GPU comparison \;\textbullet\;
20 of 20 demo images correct \;\textbullet\; No thermal throttling observed on either board
\end{tcolorbox}

\section{Overview}

The BRISC 2025 dataset provides four classes of brain MRI slice: glioma, meningioma,
pituitary tumor, and no tumor, with paired segmentation masks for the three tumor
classes. This project trains classification and segmentation baselines on that dataset,
then answers a question the dataset paper does not: what does it actually cost, in
latency, memory, power, and heat, to run these models on cheap edge hardware instead of
a workstation GPU -- and, since not all edge hardware is the same, how a CPU-only board
(the Pi) compares against a small CUDA-capable board (the Jetson Nano).

The deployed system is a two stage pipeline. A classifier decides what kind of scan it
is looking at. If the classifier is confident it is looking at a tumor, a segmenter runs
next and marks where the tumor is. Every run produces a PDF report. The whole thing has
been run and measured on real hardware -- a Raspberry Pi~4 and a Jetson Nano -- not
estimated.

\section{Models and Checkpoints}

Two classifiers and two segmenters were trained (see \texttt{train\_colab.ipynb}), so
the deployment can pick a pair by size or by accuracy. The segmenters share a ResNet34
encoder; only the decoder differs.

\begin{center}
\begin{tabular}{lllll}
\toprule
Checkpoint & Role & File size & Parameters & Reported accuracy \\
\midrule
EfficientNetB0 & Classifier & 15.59 MB & $\sim$5.3M & 99.20\% F1 \\
MobileNetV3 Large & Classifier & 16.25 MB & $\sim$5.5M & 94.47\% F1 \\
DeepLabV3+ (ResNet34) & Segmenter & 85.77 MB & $\sim$24.4M & 76.3\% mIoU \\
UNet (ResNet34) & Segmenter & 93.38 MB & $\sim$24.4M & 75.7\% mIoU \\
\bottomrule
\end{tabular}
\end{center}
\begin{center}\small Table 1: Checkpoint size, parameter count, and reported accuracy by role.\end{center}

The deployed app defaults to EfficientNetB0 and DeepLabV3+, the pair with the best
reported accuracy on each task. That default is not the fastest pair on either edge
board; see Sections~\ref{sec:pi-latency} and \ref{sec:jetson-latency} for why that trade
is still worth keeping.

\section{Inference Pipeline}

Segmentation only runs when the classifier is actually confident it is looking at a
tumor. This keeps the common case (no tumor, or a low confidence read) cheap, and avoids
drawing a mask on an image the model was never sure about in the first place.

\begin{verbatim}
1. Classify   224x224 input, ImageNet normalization -> softmax over 4 classes
2. Gate       top class is a tumor type AND confidence >= 70% ?
                    |
        yes ------- + ------- no
         |                     |
3. Segment                   skip -> verdict is "no tumor" or "uncertain"
   256x256 input,
   binary mask
         |
4. Report     verdict, confidence chart, mask overlay,
              tumor area, run metadata -> PDF
\end{verbatim}

The confidence threshold (default 70\%) is adjustable. A predicted tumor class below
threshold is reported as uncertain, not folded into either a confident tumor or a
confident no-tumor call, and segmentation is skipped either way.

\section{Edge Deployment Targets}

\begin{center}
\small
\begin{tabular}{@{}L{2.1cm}L{1cm}L{4.7cm}L{2.3cm}L{4cm}@{}}
\toprule
Board & RAM & OS & CPU & Accelerator \\
\midrule
Pi 4 Model B & 4 GB & Raspberry Pi OS Full (Bookworm, 64-bit) & 4 threads & none (CPU only) \\
Jetson Nano & 4 GB & L4T / JetPack 4.6.6 (R32.7.6, Ubuntu 18.04) & 4x Cortex-A57 & Tegra X1, 128-core Maxwell CUDA \\
\bottomrule
\end{tabular}
\end{center}
\begin{center}\small Table 2: Target hardware and OS configuration.\end{center}

64-bit Raspberry Pi OS is a hard requirement for the Pi: official PyTorch wheels only
target aarch64 Linux, and 32-bit armhf has no practical modern PyTorch build. On the
Jetson, the constraint runs the other way: the generic PyPI PyTorch wheel (CPU-only, the
same one used for the Pi) is \emph{not} built for Jetson's ARM+CUDA combination, so it
either fails to install or silently installs without GPU support. The Jetson instead
needs NVIDIA's own JetPack-matched wheel, plus torchvision built from source against it
(no prebuilt torchvision wheel exists for this combination).

\begin{verbatim}
# Pi 4 (CPU-only wheel)
pip install torch torchvision --index-url https://download.pytorch.org/whl/cpu

# Jetson Nano, JetPack 4.6.x / L4T R32.7.x (Python 3.6)
wget <NVIDIA JetPack-matched torch 1.10.0 wheel>
pip install torch-1.10.0-cp36-cp36m-linux_aarch64.whl
git clone --branch v0.11.1 https://github.com/pytorch/vision torchvision
cd torchvision && python3 setup.py install
\end{verbatim}

\section{Resource Budget}

Before anything was measured on either device, a budget was worked out on paper, then
checked against what actually happened.

\textbf{Storage.} The four checkpoints together are about 213~MB. Add the OS and Python
environment and the whole install lands around 3 to 4~GB on the Pi's 32~GB card, or a
few GB more on the Jetson once the from-source torchvision build artifacts are included.
Storage was never the binding constraint on either board.

\textbf{Memory.} The Pi's pre-deployment estimate (Table~3) put peak usage around 1.5 to
1.8~GB; the measured figure, with all four models loaded into one process at once, came
in at 676~MB peak resident memory. The Jetson runs the identical all-four-models-loaded
benchmark and measured \textbf{3023--3041~MB} peak RSS -- far higher, because the Jetson
build additionally holds CUDA context and GPU-resident weights/activations on top of the
same Python/PyTorch overhead, out of only 3964~MB total system RAM. That is close enough
to the ceiling that spawning even a small extra subprocess (e.g.\ \texttt{tegrastats})
while all four models are resident can fail with \texttt{OSError: Cannot allocate memory}
on \texttt{fork()} -- see Section~\ref{sec:jetson-ops}.

\begin{center}
\begin{tabular}{ll}
\toprule
Component & Estimated \\
\midrule
OS and services, headless & $\sim$500 MB \\
Python and the Streamlit server & $\sim$200 MB \\
PyTorch runtime import overhead & $\sim$250 MB \\
Model weights and activations & $\sim$500 MB \\
Chart and PDF rendering, transient & $\sim$100 MB \\
\bottomrule
\end{tabular}
\end{center}
\begin{center}\small Table 3: Pre-deployment memory budget estimate (Pi).\end{center}

\section{Measured Latency -- Raspberry Pi 4}
\label{sec:pi-latency}

Measured with \texttt{benchmark.py}: 2 untimed warmup runs followed by 10 timed runs per
stage, on a single 512x512 test slice, all four checkpoints loaded at once.

\begin{center}
\begin{tabular}{lrrrrr}
\toprule
Model & Mean & Median & P95 & Min & Max \\
\midrule
Classifier: EfficientNetB0 & 282.0 & 280.1 & 308.7 & 274.8 & 308.7 \\
Classifier: MobileNetV3 & 189.3 & 187.2 & 207.5 & 186.2 & 207.5 \\
Segmenter: DeepLabV3+ & 1007.6 & 1004.5 & 1045.7 & 992.6 & 1045.7 \\
Segmenter: UNet & 926.5 & 924.3 & 938.4 & 918.6 & 938.4 \\
\bottomrule
\end{tabular}
\end{center}
\begin{center}\small Table 4: Per-model latency on the Pi 4, all figures in milliseconds.\end{center}

\begin{center}
\begin{tabular}{lrrrrr}
\toprule
Combination & Mean & Median & P95 & Min & Max \\
\midrule
EfficientNetB0 + DeepLabV3+ (default) & 1295.5 & 1292.7 & 1325.5 & 1288.7 & 1325.5 \\
EfficientNetB0 + UNet & 1203.9 & 1203.1 & 1223.0 & 1194.1 & 1223.0 \\
MobileNetV3 + DeepLabV3+ & 1196.3 & 1197.8 & 1200.5 & 1188.5 & 1200.5 \\
MobileNetV3 + UNet & 1114.6 & 1114.8 & 1125.0 & 1102.2 & 1125.0 \\
\bottomrule
\end{tabular}
\end{center}
\begin{center}\small Table 5: End-to-end pipeline latency on the Pi 4, all figures in milliseconds.\end{center}

A full analysis, upload to result, lands around 1.1 to 1.3 seconds on the Pi.

\section{Measured Latency -- Jetson Nano}
\label{sec:jetson-latency}

Measured with \texttt{benchmark\_jetson.py} on the same methodology as the Pi (5 untimed
warmup runs, 20 timed runs per stage, all four checkpoints loaded at once), on CUDA
(Tegra X1 GPU, \texttt{torch.cuda.synchronize()} around every timed call since GPU kernel
launches are asynchronous).

\begin{center}
\begin{tabular}{lrrrrr}
\toprule
Model & Mean & Median & P95 & Min & Max \\
\midrule
Classifier: MobileNetV3 & 40.7 & 39.8 & 48.2 & 34.3 & 48.2 \\
Classifier: EfficientNetB0 & 51.4 & 50.5 & 59.3 & 46.2 & 59.3 \\
Segmenter: UNet & 100.7 & 100.3 & 105.1 & 99.2 & 105.1 \\
Segmenter: DeepLabV3+ & 111.7 & 111.6 & 113.7 & 111.1 & 113.7 \\
\bottomrule
\end{tabular}
\end{center}
\begin{center}\small Table 6: Per-model latency on the Jetson Nano (CUDA), all figures in milliseconds.\end{center}

\begin{center}
\begin{tabular}{lrrrrr}
\toprule
Combination & Mean & Median & P95 & Min & Max \\
\midrule
MobileNetV3 + UNet & 141.1 & 141.2 & 143.5 & 137.1 & 143.5 \\
EfficientNetB0 + UNet & 151.7 & 151.8 & 158.0 & 146.6 & 158.0 \\
MobileNetV3 + DeepLabV3+ & 151.8 & 152.0 & 154.8 & 149.0 & 154.8 \\
EfficientNetB0 + DeepLabV3+ (default) & 162.9 & 163.6 & 167.0 & 158.9 & 167.0 \\
\bottomrule
\end{tabular}
\end{center}
\begin{center}\small Table 7: End-to-end pipeline latency on the Jetson Nano, all figures in milliseconds.\end{center}

A full analysis on the Jetson lands around 140--165~ms end to end -- roughly
\textbf{8--9x faster than the Pi} on the same default model pair (162.9~ms vs
1295.5~ms), for a device in the same rough price and power class, because the Nano's
CUDA GPU can actually run these convolution-heavy models in parallel instead of walking
them through 4 CPU threads. Peak RSS with all four models loaded was
\textbf{3023--3041~MB} (see Section~5), against 676~MB on the Pi.

\section{Jetson Power and Thermal Behavior}
\label{sec:jetson-power}

Captured with \texttt{tegrastats} running concurrently in a separate SSH session while
\texttt{benchmark\_jetson.py} ran the full latency sweep above.

\begin{center}
\begin{tabular}{ll}
\toprule
Metric & Value \\
\midrule
Peak instantaneous power (\texttt{POM\_5V\_IN}) & 7309 mW ($\sim$7.3 W), during a GR3D\_FREQ 99\% segmenter pass \\
Peak GPU power (\texttt{POM\_5V\_GPU}) & 4363 mW ($\sim$4.4 W), same spike \\
Steady running-average power & $\sim$2.4--2.7 W \\
Idle baseline power & $\sim$0.9--1.5 W \\
Peak thermal (AO-therm) & 48.5\degree C -- no throttling \\
CPU clock under load & up to 1479 MHz sustained \\
\bottomrule
\end{tabular}
\end{center}
\begin{center}\small Table 8: Jetson Nano power draw and thermal peaks over the full benchmark run.\end{center}

The Jetson's thermal zones (AO, CPU, GPU, PLL, PMIC, and an iwlwifi radio zone) all
stayed well clear of any throttle point through the run -- the highest reading observed
anywhere was 48.5\degree C, versus a typical throttle point in the 90s for this SoC.
Power draw tracks GPU utilization tightly: every segmenter pass (\texttt{GR3D\_FREQ}
99\%) produces a sharp instantaneous spike toward 5--7~W, while idle and classifier-only
periods sit under 2~W.

\section{Pi versus Jetson versus T4 GPU}

The same four combinations were run with the same methodology on a Colab T4 GPU
(\texttt{colab\_gpu\_benchmark.ipynb}), with \texttt{torch.cuda.synchronize()} around
every timed call. The Jetson numbers use the clean (uncontended) run from
Section~\ref{sec:jetson-latency}.

\begin{center}
\begin{tabular}{lrrrr}
\toprule
Model & Pi 4 (CPU) & Jetson Nano (CUDA) & T4 (CUDA) & Pi $\to$ Jetson speedup \\
\midrule
Classifier: EfficientNetB0 & 282.0 ms & 51.4 ms & 8.76 ms & 5.5x \\
Classifier: MobileNetV3 & 189.3 ms & 40.7 ms & 6.10 ms & 4.7x \\
Segmenter: DeepLabV3+ & 1007.6 ms & 111.7 ms & 9.76 ms & 9.0x \\
Segmenter: UNet & 926.5 ms & 100.7 ms & 6.51 ms & 9.2x \\
\bottomrule
\end{tabular}
\end{center}
\begin{center}\small Table 9: Mean per-model latency across all three platforms.\end{center}

The Jetson Nano sits where its hardware suggests it should: a genuine CUDA GPU (even a
small, several-generations-old Maxwell part) closes most of the gap to a datacenter T4 on
the compute-heavy segmenters, while still being $\sim$10--30x slower than the T4 itself.
Segmentation again gains more from having any GPU at all (9.0--9.2x Pi-to-Jetson) than
classification does (4.7--5.5x), consistent with the Pi-vs-T4 finding in the original
report: the heavier ResNet34-encoder segmenters have more to parallelize than the
lighter classifiers.

Read this three-way comparison as directional, not exact -- as with the original Pi/T4
comparison, the three runs used different test images and run counts.

\section{Correctness Check}

Speed is only useful if the answers are right. All 20 demo images (5 per class) were run
through the default pair on the T4 GPU; every prediction matched its known class, every
tumor image produced a segmentation with a plausible tumor area, and every no-tumor image
correctly skipped segmentation.

\begin{center}
\begin{tabular}{llll}
\toprule
Class & Predicted & Confidence range & Result \\
\midrule
Glioma (5 images) & glioma & 92.1\% to 100.0\% & segmented \\
Meningioma (5 images) & meningioma & 99.6\% to 100.0\% & segmented \\
Pituitary (5 images) & pituitary & 99.9\% to 100.0\% & segmented \\
No tumor (5 images) & no\_tumor & 100.0\% & skipped \\
\bottomrule
\end{tabular}
\end{center}
\begin{center}\small Table 10: Prediction outcome on the 20-image demo set (T4 GPU, default pair).\end{center}

\section{Wireless Scan Pipeline}

The deployment also covers getting a raw scan onto an edge board in the first place.
This half runs on a separate local computer. Two send targets are supported --
\texttt{send\_to\_pi.py} and \texttt{send\_to\_jetson.py} -- sharing the same SSH/SCP
plumbing (\texttt{ssh\_utils.py}) so the pipeline works identically against either board.

\begin{verbatim}
.nii / .nii.gz scan
  | nii_to_slices.py, local computer
  v
PNG slices, anonymized sequential filenames
  | send_to_pi.py or send_to_jetson.py, SSH / SCP
  v
case_report.py / case_report_jetson.py runs on the target board
  | classifies every slice, ranks by confidence, segments only
  | the single best view, writes one consolidated PDF
  v
PDF report
  | SCP back to the local computer
  v
Downloaded report
\end{verbatim}

\texttt{nii\_web\_app.py} wraps the upload, conversion, preview, and send steps (to
either board) in a browser page so the local side needs no terminal. It requires SSH key
based authentication, because a web app cannot answer an interactive password prompt; it
fails fast with a clear message (including a host-key verification failure, the first
time a given machine talks to a given board) instead of hanging if auth has not been set
up yet.

\section{Operational Notes}
\label{sec:jetson-ops}

A short record of the real issues hit while wiring this up on actual hardware, kept here
so the next setup does not repeat the same debugging.

\subsection*{Raspberry Pi}
\begin{itemize}[leftmargin=1.2em]
\item PyTorch 2.6 changed \texttt{torch.load}'s default. Loading a checkpoint that is
  not a bare tensor now fails unless \texttt{weights\_only=False} is passed explicitly.
\item Git LFS was not installed on the Pi at first; checkpoints pulled as 133 byte
  pointer stubs instead of real weights.
\item Git Bash on Windows rewrites POSIX-looking paths in remote SSH command strings;
  \texttt{MSYS\_NO\_PATHCONV=1} on the subprocess environment stops it.
\end{itemize}

\subsection*{Jetson Nano}
\begin{itemize}[leftmargin=1.2em]
\item The generic PyPI PyTorch wheel is not built for Jetson's ARM+CUDA combination and
  must be replaced with NVIDIA's JetPack-matched wheel; torchvision has no matching
  prebuilt wheel at all and must be built from source against it.
\item Older JetPack images ship Python 3.6 with an old pip/setuptools that has no
  prebuilt wheel for numpy, Pillow, or matplotlib on aarch64 -- installing these via
  \texttt{apt} (\texttt{python3-numpy}, \texttt{python3-pil}, \texttt{python3-matplotlib})
  and inheriting them into a \texttt{--system-site-packages} venv avoids a failed
  from-source compile.
\item \texttt{torch.load}'s \texttt{weights\_only} keyword does not exist on torch
  1.10.0 (the JetPack 4.6.x wheel) -- it was added in torch 1.13. Passing it raises
  \texttt{TypeError}; omitting it is behavior-preserving, since full unpickling was
  already torch 1.10's only mode.
\item \texttt{subprocess.Popen(..., text=True)} is Python 3.7+ only; Jetson's stock
  Python 3.6 needs \texttt{universal\_newlines=True} instead, which behaves identically.
\item With all four models resident (peak RSS $\sim$3~GB out of 3964~MB total), spawning
  even a small subprocess like \texttt{tegrastats} can fail \texttt{fork()} with
  \texttt{OSError: [Errno 12] Cannot allocate memory} rather than the more familiar
  \texttt{FileNotFoundError} case (binary absent, e.g.\ off-Jetson) -- both are now caught
  the same way, degrading to "power draw not captured" instead of crashing the benchmark.
\item Running \texttt{tegrastats} concurrently in a second SSH session while the
  benchmark runs (rather than having the benchmark spawn it inline) avoids that memory
  contention and also avoids the two processes competing for CPU/GPU time, which was
  observed to roughly double tail latency (p95) versus an uncontended run.
\end{itemize}

\section{Limitations}

\textbf{Not a diagnostic device.} This is a research and demonstration system. Every
report it produces says so. Findings should always be reviewed by a qualified
radiologist before they inform any clinical decision.

\textbf{Domain mismatch is a real failure mode, not a hypothetical one.} The models were
trained on BRISC: contrast-enhanced T1-weighted slices with the skull present. Running a
scan from a different acquisition pipeline produced plausible-looking but untrustworthy
classifications during testing. The pipeline does not know when it has been handed
out-of-distribution data; it will answer confidently either way.

\textbf{Benchmark comparability.} The Pi, Jetson, and T4 timing runs were not a
perfectly controlled A/B: different single test images, different run counts, and
different hardware generations. Treat the speedup figures in Section~9 as
representative, not as a certified benchmark result.

\vspace{1em}
\noindent\rule{\linewidth}{0.4pt}
\small\color{midgray} BRISC 2025 -- inference and edge deployment report. Generated from measured runs, no simulated figures.

\end{document}
